\documentclass{article}
\input{../../../lib.sty}

\title{\texttt{fn sample\_poisson}}
\author{OpenDP}

\begin{document}
\maketitle

\contrib Proves soundness of \texttt{fn sample\_poisson} in \texttt{mod.rs}.
This proof reduces the general Poisson sampler to repeated calls to
\texttt{sample\_poisson\_0\_1} on equal sub-means in $[0,1)$, then sums the results.

\section{Hoare Triple}

\subsection*{Precondition}
\subsubsection*{Compiler-verified}
\begin{itemize}
    \item Argument \texttt{lambda} is of type \texttt{RBig}, a bignum rational
\end{itemize}

\subsubsection*{User-verified}
$\texttt{lambda} \ge 0$

\subsection*{Implementation}
\lstinputlisting[language=Python,firstline=2,escapechar=|]{./pseudocode/sample_poisson.py}

\subsection*{Postcondition}
\begin{theorem}
\label{postcondition}
For any setting of the input parameter \texttt{lambda} such that the given preconditions hold, \\
\texttt{sample\_poisson} either returns \texttt{Err(e)} due to a lack of system entropy,
or \texttt{Ok(out)}, where \texttt{out} is distributed as $Poisson(\texttt{lambda})$.
\end{theorem}

\section{Proof}

\begin{lemma}
\label{err-e}
\texttt{sample\_poisson} only returns \texttt{Err(e)} when there is a lack of system entropy.
\end{lemma}

\begin{proof}
If \texttt{lambda} $= 0$, the function returns \texttt{Ok(0)} immediately.

Otherwise, the only fallible call is line \ref{line:subcall}, which invokes
\rustdoc{traits/samplers/poisson/fn}{sample\_poisson\_0\_1}.
By line \ref{line:pieces}, \texttt{pieces} $= \lfloor \texttt{lambda} \rfloor + 1$, so
\texttt{pieces} is a positive integer. Therefore line \ref{line:piece} defines
\[
\texttt{piece} = \frac{\texttt{lambda}}{\texttt{pieces}},
\]
which satisfies $0 \le \texttt{piece} < 1$ because $\texttt{pieces} > \texttt{lambda}$ whenever
$\texttt{lambda} > 0$. Hence the preconditions of \texttt{sample\_poisson\_0\_1} hold.

By the definition of \texttt{sample\_poisson\_0\_1}, any error it returns is due to a lack of
system entropy. Therefore \texttt{sample\_poisson} only returns \texttt{Err(e)} when there is
a lack of system entropy.
\end{proof}

We now condition on not returning an error.

\begin{lemma}
\label{piece-law}
Each invocation of \texttt{sample\_poisson\_0\_1(piece)} on line \ref{line:subcall}
returns an independent random variable distributed as $Poisson(\texttt{piece})$.
\end{lemma}

\begin{proof}
By the proof of Lemma~\ref{err-e}, the preconditions of
\texttt{sample\_poisson\_0\_1} hold for \texttt{piece}. Therefore each call returns a
$Poisson(\texttt{piece})$ random variable by its postcondition. Independence follows from the
fact that the calls are made sequentially with fresh randomness.
\end{proof}

\begin{lemma}
\label{sum-poisson}
If $X_1,\ldots,X_m$ are independent and each $X_i \sim Poisson(\mu)$, then
\[
X_1 + \cdots + X_m \sim Poisson(m\mu).
\]
\end{lemma}

\begin{proof}
The probability generating function of a $Poisson(\mu)$ random variable is
\[
G_X(z) = \mathbb{E}[z^X] = \sum_{k=0}^\infty e^{-\mu}\frac{\mu^k}{k!}z^k = e^{\mu(z-1)}.
\]
Therefore, if $X_1,\ldots,X_m$ are independent and each distributed as $Poisson(\mu)$,
the probability generating function of their sum is
\[
\mathbb{E}[z^{X_1+\cdots+X_m}]
=
\prod_{i=1}^m \mathbb{E}[z^{X_i}]
=
\left(e^{\mu(z-1)}\right)^m
=
e^{m\mu(z-1)},
\]
which is the probability generating function of $Poisson(m\mu)$.
\end{proof}

\begin{lemma}
\label{ok-out}
If the outcome of \texttt{sample\_poisson} is \texttt{Ok(out)}, then \texttt{out} is distributed as $Poisson(\texttt{lambda})$.
\end{lemma}

\begin{proof}
If \texttt{lambda} $= 0$, the function returns $0$, which is $Poisson(0)$.

Assume now that \texttt{lambda} $> 0$. Let
\[
m := \lfloor \texttt{lambda} \rfloor + 1
\quad\text{and}\quad
\mu := \frac{\texttt{lambda}}{m}.
\]
By Lemma~\ref{piece-law}, each iteration of the loop on line \ref{line:subcall} adds an
independent $Poisson(\mu)$ random variable to \texttt{total}. Since the loop runs exactly $m$
times, \texttt{out} is the sum of $m$ independent $Poisson(\mu)$ random variables.

By Lemma~\ref{sum-poisson},
\[
\texttt{out} \sim Poisson(m\mu) = Poisson(\texttt{lambda}).
\]
Thus \texttt{out} has the claimed distribution.
\end{proof}

\begin{proof}[Proof of Theorem~\ref{postcondition}]
Holds by Lemma~\ref{err-e} and Lemma~\ref{ok-out}.
\end{proof}

\bibliographystyle{alpha}
\bibliography{mod}
\end{document}